% =============================================================
\section{简介}
\label{sec:intro}
% =============================================================

与可用六自由度位姿描述的刚体不同，毛巾等柔性织物在抓取与搬运中容易发生褶皱、卷边、层间滑移与自遮挡，难以依靠解析动力学模型预测其形变过程。常规二指平行夹爪仅具备单一开合自由度，在处理多层织物捏持、折边对齐与翻面动作时容易漏夹下层织物或拉乱已成形折边。本项目以两台松灵 AgileX Nero 七自由度机械臂与两套灵心巧手 LinkerHand L20 多指灵巧手作为执行机构，配合单台 Intel RealSense L515 激光雷达相机提供全局对齐 RGB-D 观测，研究 46 自由度双臂灵巧手系统的遥操作示教学习与自主连续折布控制方法。

课程项目要求系统从工作垫上任意初始位置与平面偏航角的平铺矩形毛巾出发，在不依赖人工中途干预与辅助工装的条件下，自主完成连续两次对折，将毛巾折叠为四层叠放结构并留置于黑色工作垫中央区域。此外，系统需在有限示教样本规模下具备单侧抓空后的在线补抓恢复能力，并在多进程软硬件架构下保证双臂多指协调运动的实时性与物理安全。

围绕上述要求，本项目完成了双臂相机与虚拟工具中心点（TCP）标定、视觉遥操作数据采集、模仿学习策略训练以及代码生成智能体的闭环集成与实机测试。在遥操作与数据集构建方面，团队基于 Pico 4 头显的手部骨骼追踪搭建了双臂增量遥操作平台，结合 Pinocchio 解析雅可比逆运动学与对指距离自适应重定向算法，分别录制了一次对折与二次对折双臂示教轨迹各 60 条，统一导出为 LeRobot v3.0 格式数据集。在模仿学习方面，训练了面向 46 维联合动作空间的动作分块变换器（Action Chunking with Transformers, ACT）~\cite{zhao2023learning}基线策略，并实现了采用冻结 DINOv3 ViT-S/16~\cite{simeoni2025dinov3}预训练视觉主干与标准化本体状态高斯噪声正则的表征改进策略。在实机部署阶段，实验对比表明单阶段 60 条示教轨迹训练的神经网络在大范围随机初始摆放条件下容易发生抓取偏移，而基于三维视觉反投影的代码生成智能体在一次对折与居中摆正任务上对初始位姿变化更为稳健。最终系统采用双阶段架构，先由基于模型上下文协议（Model Context Protocol, MCP）与持久工作进程的双臂折布智能体完成一次对折并将双层毛巾翻面移至垫面中央，再经双臂关节回位与首帧手型对齐后交由双臂 ACT 策略完成二次对折。


% =============================================================
\section{问题描述}
\label{sec:problem}
% =============================================================

本项目的控制对象是单目 RGB-D 视觉引导下的双臂多指灵巧手连续折布系统。相比平行夹爪桌面操作，该任务在运动学求解、坐标标定、分阶段控制方式选择与跨模块衔接上存在四方面具体约束。

第一，执行机构自由度高且存在人机形态差异。单侧 Nero 机械臂含 7 个旋转关节，单侧 LinkerHand L20 灵巧手含 16 个独立受控主动槽位与 5 个机械耦合被动指节，左右两侧合计构成 46 维连续状态与动作空间。人手骨骼与机器人灵巧手的指骨长度及关节轴向并不相同，直接复刻人手关节角会导致拇指与其余手指无法闭合接触，因此需要在每个控制周期内同步求解带限位约束的七轴连续逆运动学与五指运动学重定向。同时，多指灵巧手不存在平行夹爪那样可直接量测的对称几何中心，相机光学系、左右机械臂基座系与手指实际拈布点之间的刚体变换必须通过手眼与虚拟工具联合标定确定。

第二，连续两次对折在前后两个阶段面临不同的主要矛盾。第一次对折时，单层毛巾随机平铺于桌面，中心位置与偏航角在试验间变化幅度较大，核心需求是在大范围平面内定位待折角点、使指尖捏合平面顺应布边切向，并在偶发滑脱后单独补抓滞留角点。仅凭 60 条遥操作轨迹很难密集覆盖二维位置与偏航角的组合空间，落点数厘米的预测偏差便会造成抓空或折边严重错位。进入第二次对折后，毛巾已被折为居中的双层长方形，此时双臂需要在折边两端同时捏住上下两层织物并协调完成抬升、弧线翻折与释放，动作好坏取决于多指捏持配合与双臂同步轨迹，难以仅靠少量人工设定的离散几何航点复现自然协调的折叠过程。

基于上述分析，本项目设计并实现了一套分层耦合的双臂灵巧操作软硬件系统。系统在功能上涵盖双臂同步遥操作录制、九自由度相机与虚拟 TCP 联合标定、46 维 ACT 策略训练与冻结视觉主干改进，以及代码生成式任务规划与跨阶段状态对齐执行。在工程运行上，底层 CAN 总线控制环境与 GPU 深度学习推理环境通过进程隔离解耦，所有机械臂目标在执行前均经过关节限位与笛卡尔工作空间立方体校验，并配备独立于上层模型会话的人工位置保持旁路以确保实机安全。
